\section{Background}
\label{sec:background}

\subsection{Extensive-form games}

A two-player zero-sum extensive-form game is defined by a game tree where
each node is either a chance node (nature acts), a decision node (a player
acts), or a terminal node (payoffs are realized). HUNL is a two-player
zero-sum game in extensive form with imperfect information.

Formally: a game
$\Gamma = (N, H, Z, A, u, \sigma_c, \mathcal{I})$ where $N = \{1, 2\}$ is the
player set, $H$ is the set of histories, $Z \subseteq H$ is the set of
terminal nodes, $A(h)$ is the action set at history $h$,
$u_i\colon Z \to \mathbb{R}$ is the utility function for player $i$,
$\sigma_c$ is the chance strategy, and $\mathcal{I}$ is the information
partition.

\subsection{Nash equilibrium and exploitability}

A \emph{Nash equilibrium} is a strategy profile
$(\sigma_1^\star, \sigma_2^\star)$ such that neither player can improve their
expected utility by unilaterally deviating:
\[
  u_i(\sigma_i^\star, \sigma_{-i}^\star)
    \;\geq\; u_i(\sigma_i, \sigma_{-i}^\star)
    \quad \forall \sigma_i, \forall i \in N.
\]
In two-player zero-sum games, Nash equilibria always exist and all Nash
equilibria yield the same expected payoffs (von Neumann's minimax theorem).

\emph{Exploitability} of a strategy $\sigma_i$ is the maximum gain available
to a best-responding opponent:
\[
  \mathrm{expl}(\sigma_i)
    \;=\; \max_{\sigma_{-i}} u_{-i}(\sigma_i, \sigma_{-i}) - v^\star,
\]
where $v^\star$ is the game value. Exploitability of $0$ means $\sigma_i$ is
a Nash strategy. We use exploitability as our primary convergence metric.

\subsection{Counterfactual regret minimization (CFR)}

CFR \citep{zinkevich2007} solves for Nash equilibria in extensive-form games
by iteratively minimizing regret at each information set. At each iteration
$t$, for each infoset $I$ and action $a$:
\[
  r^t(I, a)
    \;=\; \sum_{h \in I} \pi_{-i}^{\sigma^t}(h)
          \left[v_i(\sigma^t_{I \to a}, h) - v_i(\sigma^t, h)\right],
\]
with cumulative regret $R^T(I, a) = \sum_{t=1}^T r^t(I, a)$ and the
regret-matching strategy update
\[
  \sigma^{T+1}(I, a)
    \;=\; \frac{R^T_+(I, a)}{\sum_{a'} R^T_+(I, a')},
\]
where $R^T_+(I, a) = \max(R^T(I, a), 0)$. The average strategy
$\bar{\sigma}^T$ has exploitability bounded by $O(1/\sqrt{T})$.

\subsection{CFR$^+$}

CFR$^+$ \citep{tammelin2014} modifies the regret accumulation step:
\[
  R^{T+1}(I, a)
    \;=\; \max\!\left(R^T(I, a) + r^T(I, a), \; 0\right).
\]
Negative cumulative regrets are floored to zero at each iteration rather
than at strategy computation time. This seemingly small change yields
dramatically faster empirical convergence---in practice, CFR$^+$ typically
requires an order of magnitude fewer iterations than vanilla CFR to reach
the same exploitability level.

\subsection{Game abstraction}

Full HUNL has approximately $10^{18}$ distinct infosets---computationally
intractable for exact solving. Practical solvers work within an abstracted
game.

\paragraph{Card abstraction} groups strategically similar hands into buckets,
replacing exact hand identities with equity-percentile estimates. The
abstracted game has $K$ card buckets instead of the full hand space.

\paragraph{Action abstraction} discretizes the continuous bet-size space
into a finite set of sizes, reducing the branching factor at each decision
point from infinite to a small fixed number.
